\section{Related Works}

% \daniel{consider adding paragraph heading so that the reader can expect what each paragraph is about.}

\paragraph{Summarization evaluation.} PLS research often introduces either datasets~\cite{goldsack2022making,Crossley2021TheCE,pu-etal-2024-scinews,Manor2019PlainES}, methods~\cite{Guo2022RetrievalAO,August2022GeneratingSD,luo2022readability,ji-etal-2024-rag,flores-etal-2023-medical}, or both~\cite{Guo2020AutomatedLL,Zaman2020HTSSAN,chandrasekaran2020overview}. The majority of prior work use a combination of readability metrics, such as Flesch Reading Ease~\cite{Flesch1943MarksOR} or the Gunning-Fog Index~\cite{Gunning1968TheTO} to validate the readability of their dataset or generations. Readability metrics are typically reported in conjunction with more general summarization metrics, such as ROUGE~\cite{lin2004rouge} or BertScore~\cite{zhang2020bertscore}. General summarization evaluation is a well-studied area, with ongoing work analyzing both the efficacy of summarization metrics~\cite{fabbri2021summeval,khashabi2022genie,goyal2022news} and designing metrics that better align with human judgments~\cite{liu-etal-2023-towards-interpretable,Liu2022RevisitingTG}. \textcite{guo-etal-2024-appls} analyzed how perturbations in plain language summaries affect results of general summarization metrics. In this chapter, we focus on readability metrics, rather than general summarization metrics, with the goal of understanding how well readability metrics measure readability for PLS.

\paragraph{Readability Metrics.} While readability metrics are well studied in fields such as education~\cite{gray1935makes, dubay2004principles, texts_2022} and linguistics~\cite{Pires02102017}, there is little work studying how well these metrics perform for the task of PLS. 
Most traditional metrics were not designed specifically for PLS, or even for evaluation in Computer Science. The most common origin of traditional metrics is the need to assess the readability of K-12 school texts~\cite{dale1948formula, Coleman1975ACR}. %Flesch Reading Ease was developed for the U.S. Navy to asses the difficulty of technical manuals~\cite{Flesch1943MarksOR}. 
Linsear Write was introduced in the book, \textit{Gobbledygook has gotta go}, published by the US Department of the Interior for the purposes of measuring the complexity of government communications~\cite{o1966gobbledygook}.
% Coleman-Liau was published for the purposes of rating the readability of public school textbooks on a large scale, with the specific benefit of only needing a an optical scanning device to run the calculations~\cite{Coleman1975ACR}. %In contrast, the majority of PLS literature focuses on generating summaries of scientific papers for a general audience, in which ``general audience'' is defined as an adult who is not an expert in the paper's domain~\todocite. The inherent difference in audience (K-12 vs general adult population) and intent (school-grade education vs scientific communication), means that these metrics may not cleanly apply to the task of PLS. In this paper, we attempt to answer the question of application of readability metrics to PLS.
As readability metrics rely primarily on lexical features~\cite{508dd5a7-e076-39ee-bf49-06f80e20c28d}, prior work has offered criticism of readability metrics, showing that they can be easily manipulated to provide better scores with changes that do not substantially improve the readability of summaries~\cite{Tanprasert2021FleschKincaidIN}. Other work has looked at which linguistic attributes are correlated with readability metrics~\cite{tajner2012WhatCR}. To the best of our knowledge, our work is the first to measure the correlation of readability metrics with human readability judgments.

\paragraph{LMs as Evaluators.} Recent advances in LMs have shown that they are capable of reasoning over complex language~\cite{Brown2020LanguageMA, Wei2022ChainOT, Yang2024DoLL}. LMs have been shown to be effective evaluators in other natural language tasks~\cite{li-etal-2025-exploring-reliability,10.1145/3640457.3688075,nedelchev-etal-2020-language,liu2023geval}, including related summarization tasks~\cite{song-etal-2024-finesure}. Given this success in prior work, we hypothesize that LMs are capable of evaluating the readability of plain language summaries. In particular, we hypothesize that LMs can reason over more complex attributes of readability, such as the background required or whether technical concepts are explained. 
